\pdfbookmark[0]{封面}{titlepage}
\begin{titlepage}
	\rightline{\includegraphics[scale=0.6]{logo}}
	\begin{center}
		\vspace*{4cm}
		\textbf{\heiti \zihao{1} 防\enspace 灾\enspace 科\enspace 技\enspace 学\enspace 院}
		
		\vspace*{1cm}
		\textbf{\kaishu \zihao{-0} 毕\enspace 业\enspace 设\enspace 计}
		
		\vspace*{6cm}
	\begin{tabular}{cp{4cm}cp{4cm}}
		\textbf{\zihao{4} 题\qquad 目}&\multicolumn{3}{c}{\zihao{4} 基于GUI的误差理论与数据处理}\\\cline{2-4}
		&&&\\
		\textbf{\zihao{4} 学生姓名}&\multicolumn{1}{c}{\zihao{4} 艾江南}&\textbf{\zihao{4} 学\qquad 号}&\multicolumn{1}{c}{\zihao{4} 135032105}\\\cline{2-2}\cline{4-4}
		&&&\\
		\textbf{\zihao{4} 班\qquad 级}&\multicolumn{1}{c}{\zihao{4} 1350321}&\textbf{\zihao{4} 专\qquad 业}&\multicolumn{1}{c}{\zihao{4} 测控技术与仪器}\\\cline{2-2}\cline{4-4}
		&&&\\
		\textbf{\zihao{4} 系\qquad 别}&\multicolumn{3}{c}{\zihao{4} 防灾仪器系}\\\cline{2-4}
		&&&\\
		\textbf{\zihao{4} 开题时间}&\multicolumn{1}{c}{\zihao{4} 2016年12月26日}&\textbf{\zihao{4} 答辩时间}&\multicolumn{1}{c}{\zihao{4} 2017年x月x日}\\\cline{2-2}\cline{4-4}
		&&&\\
		\textbf{\zihao{4} 指导教师}&\multicolumn{1}{c}{\zihao{4} 姚振静}&\textbf{\zihao{4} 职\qquad 称}&\multicolumn{1}{c}{\zihao{4} 副教授}\\\cline{2-2}\cline{4-4}
	\end{tabular}
	\end{center}
\end{titlepage}
\newpage
\thispagestyle{empty}
\pdfbookmark[0]{摘要}{center}
\begin{center}
	\textbf{\heiti \zihao{3} 基于GUI的误差理论与数据处理}
	
	\vspace{1em}
	作\qquad 者\quad 艾江南
	
	指导老师\quad 姚振静
\end{center}
\vspace{2em}
\paragraph{摘要}{\kaishu \zihao{5} 传统误差数据处理要花费很长时间，为了使用计算机软件辅助数据处理的要求，设计了基于Matlab GUI的误差理论与数据处理系统。设计的系统以传统数据处理的算法为核心，编写Matlab应用程序，完成了测量数据基本处理，最小二乘法处理，误差的合成，测量不确定度和回归分析的Matlab算法设计，以及对应的GUI界面。整个系统具有数据的导入、处理、显示和保存的功能。软件结果表明，设计的误差理论与数据处理软件比传统的数据处理更方便直观，并且能快速得出结果，节省大量的时间。}

\vspace{2em}
\paragraph{关键词:}\textbf{\zihao{4} 误差理论；数据处理；Matlab；GUI}
\newpage
\thispagestyle{empty}
\pdfbookmark[0]{Abstract}{0}
\begin{center}
	\textbf{\heiti \zihao{3} Title:XXXXXXXXXXXXXXXXXXXXXX}
	
	\vspace{1em}
	Author\quad Ai Jiangnan
	
	Tutor\quad Yao zhenjing
\end{center}
\vspace{2em}
\paragraph{Abstract}{\zihao{5} XXXXXXXXXXXXXXXXXXXXXXXXX.}
\vspace{2em}
\paragraph{Keywords:}\textbf{\zihao{4} Keywords1;Keywords2;Keywords3}
\newpage
\pdfbookmark[0]{目录}{1}
\begin{center}
	\tableofcontents
\end{center}
\thispagestyle{empty}